\documentclass[pdflatex,sn-mathphys-num]{sn-jnl}

\usepackage{graphicx}%
\usepackage{multirow}%
\usepackage{amsmath,amssymb,amsfonts}%
\usepackage{amsthm}%
\usepackage{mathrsfs}%
\usepackage[title]{appendix}%
\usepackage{xcolor}%
\usepackage{textcomp}%
\usepackage{manyfoot}%
\usepackage{booktabs}%
\usepackage{algorithm}%
\usepackage{algorithmicx}%
\usepackage{algpseudocode}%
\usepackage{listings}%
\usepackage{lmodern}%
\usepackage{siunitx}
\usepackage{microtype}
\usepackage{derivative}
\usepackage[version=4]{mhchem}
\usepackage{cleveref}
\usepackage[spanish]{babel}

\renewcommand{\arraystretch}{1.2}
\sisetup{group-digits=true,
  group-separator={\,},
  separate-uncertainty
}

\theoremstyle{thmstyleone}%
\newtheorem{theorem}{Teorema}%
\newtheorem{proposition}[theorem]{Proposición}%

\theoremstyle{thmstyletwo}%
\newtheorem{example}{Ejemplo}%
\newtheorem{remark}{Observación}%

\theoremstyle{thmstylethree}%
\newtheorem{definition}{Definición}%

\DeclareSIUnit\angstrom{\text{\r{A}}}

\raggedbottom
\unnumbered% uncomment this for unnumbered level heads

\begin{document}

\title[Nanoestructuras de \ce{TiO2}]{Física de Nanoestructuras de \ce{TiO2} y
Biosensores: Resumen de la Charla del Prof. J. A. Calderón Córbita}

\author{\fnm{Julian L.}
\sur{Avila-Martinez}}\email{jlavilam@udistrital.edu.co}

\affil{\orgdiv{Programa Académico de Física}, \orgname{Universidad Distrital
Francisco José de Caldas}, \orgaddress{\city{Bogotá D.C.},
\country{Colombia}}}

\abstract{
Se presenta una formulación sintética y conceptual sobre la termodinámica de
formación y propiedades físicas de nanoestructuras de dióxido de titanio.
Se contrastan los mecanismos de nucleación por deposición química frente a la
anodización electroquímica. Adicionalmente, se precisan
las fases cristalográficas del material, su comportamiento como cristal fotónico
y las condiciones de frontera electromagnéticas que posibilitan la transducción
capacitiva en arreglos de biosensores.
}

\keywords{Dióxido de titanio, anodización electroquímica, cristales fotónicos,
minimización de energía, biosensores capacitivos}

\maketitle

\section{Termodinámica de la Síntesis y Auto-organización}
\label{sec:sintesis}

La generación de heteroestructuras en estado sólido exige el control preciso
del transporte de masa y energía en la interfase dieléctrico-metal.
El crecimiento mediante deposición por baño químico obedece a fuerzas de
adsorción débiles, como las interacciones de Van der Waals o puentes de
hidrógeno, requiriendo una excitación térmica del sustrato amorfo para inducir
la nucleación.

En riguroso contraste, la síntesis de nanotubos de \ce{TiO2} constituye un
mecanismo reducto-oxidativo y corrosivo impulsado netamente por un campo
eléctrico externo. Al sumergir un ánodo de titanio
en un electrolito fluorado acuoso, emerge una dinámica competitiva: la rápida
oxidación del metal de transición y la subsecuente disolución espacialmente
localizada de dicho óxido mediada por iones flúor.
La minimización del funcional de acción y la energía libre del sistema exige
que la configuración morfológica colapse de manera espontánea en una red
hexagonal compacta, optimizando la distribución espacial y maximizando el
área superficial activa. La geometría resultante
depende de la interdependencia estricta entre el potencial aplicado,
la concentración iónica y la estequiometría del solvente.

\section{Cristalografía y Respuesta Dieléctrica}
\label{sec:cristalografia}

La estabilización cristalográfica del \ce{TiO2} exhibe un acentuado
polimorfismo térmico. Las configuraciones de red evolucionan hacia estructuras
de rutilo o anatasa tras procesos de recocido, siendo el rutilo el estado base
de menor energía y mayor estabilidad mecánica frente a la ruptura dieléctrica.

En el régimen de dispersión electrónica, el \ce{TiO2} opera como un
semiconductor de banda ancha, caracterizado por una brecha de energía prohibida
comprendida entre \qty{3.0}{\electronvolt} y \qty{3.3}{\electronvolt}.
Esta asimetría en los estados de conducción induce una alta
transmitancia en el espectro visible (longitudes de onda superiores a
\qty{510}{\nano\meter}) mientras confina la absorción electromagnética al
régimen ultravioleta.

La morfología tubular introduce una perturbación nanométrica periódica en la
constante dieléctrica del medio (alternando \ce{TiO2} y vacío). Esta asimetría
espacial actúa como una red de difracción, filtrando selectivamente los modos
electromagnéticos permitidos y otorgando a la matriz un carácter de cristal
fotónico con fuerte respuesta iridiscente.
Adicionalmente, se puede modificar el comportamiento puramente diamagnético del
óxido al dopar la matriz cristalográfica con impurezas de cobalto. Al sustituir
sitios atómicos de titanio, el cobalto aporta espines desapareados en la capa
$3d$, introduciendo momentos magnéticos acoplables a campos externos.

\section{Transducción Funcional en Biosensores}
\label{sec:biosensores}

La configuración topológica de los nanotubos provee una plataforma idónea para
la inmovilización de secuencias de oligonucleótidos. El diseño
de estas cadenas de ADN con alta afinidad selectiva hacia analitos específicos
(e.g., ketamina) induce un estado ligado dictado por el decaimiento abrupto de
la energía libre de Gibbs. La captura electroquímica del
analito desencadena una transición electrónica espontánea, inyectando
portadores de carga a través de grupos fosfato y nitrilo hacia la red
semiconductora subyacente. Este desbalance en la densidad de
estados superficial se manifiesta de forma determinista como una perturbación
puramente capacitiva, consolidando al sistema como un transductor
optoelectrónico de alta resolución geométrica.

\end{document}
